\clearpage
\fancyhead[L]{\small GAUSS AN VON LINDENAU.}
\hypertarget{native-316-lower}{}
\begin{center}
\textbf{GAUSS AN VON LINDENAU.}

\par\vspace{4pt}\hrule\vspace{4pt}
\fontsize{10}{12}\selectfont Monatliche Correspondenz. Band XIX. Seite 501..503. 1809. Mai.
\par\vspace{4pt}\hrule\vspace{4pt}
\end{center}

\begin{flushright}Göttingen, am 20. Mai 1809.\end{flushright}

\par\indent In der Anlage überschicke ich Ihnen die Resultate meiner zum Theil bereits seit einiger Zeit beendigten Rechnungen über die Verbesserung der Elemente der neuen Planeten. Bei der Ceres werden Sie einige Unterschiede zwischen meinen und \textsc{Santini}'s Resultaten (Monatl. Corresp. Febr. S. 191) wahrnehmen. Wahrscheinlich haben sie ihren Grund zum Theil darin, dass \textsc{Santini} bei der Vergleichung der \emph{einzelnen} Beobachtungen mit den XI. Elementen nicht auf die Störungen Rücksicht genommen hat, welches ich nicht billigen kann, weil diese Störungen 22 Tage hindurch keinesweges einen gleichen Einfluss auf die geocentrischen Örter haben, daher die mittlere Differenz nicht als die wahre angenommen werden darf. Auch bei der nachherigen Berechnung der Störungen scheint Herr \textsc{Santini} einige Fehler begangen zu haben; ich finde nach meinen Tafeln Störung der Länge $+6'\;6''4$, im Radius Vector $+0{,}001570$. Sie werden bemerken, dass diese Zahlen mit den \textsc{Santini}schen fast keine Ähnlichkeit haben; \SourceLeaf{317}die Störung der Breite, die sich freilich mit der \textsc{Santini}'schen nicht vergleichen lässt, weil ich nach vollständigern Gleichungen gerechnet habe, ist $+56''8$. — Auf beiliegendem Blatte schicke ich Ihnen zugleich die Ephemeriden der Ceres, Juno und Vesta für ihre nächste Sichtbarkeit. Die für die Juno hat Herr Professor \textsc{Harding} berechnet nach den schon früher in der Monatl. Corresp. abgedruckten Elementen. Die Ephemeride der Vesta nach den IV Elementen ist von Herrn Doctor \textsc{Schumacher}, welcher sich jetzt bei uns aufhält. Die für die Ceres habe ich selbst übernommen. [\textsc{Von Lindenau} hat diese Ephemeride, welche mit viertägigen Zwischenräumen für Mitternacht in Göttingen von 1809 Juli 18 bis 1810 März 16 die geocentrische gerade Aufsteigung und Abweichung auf Minuten den Logarithmus des Abstandes in vier Decimalen angibt in dem Aufsatze ,,Fortgesetzte Nachrichten u.\,s.\,f.“ Mon. Corr. 1809 Mai-Heft Seite 504 bis 515 aufgenommen.] Es fehlt also nur noch die Pallas. Leider habe ich von dieser noch keine brauchbaren Beobachtungen aus dem vorigen Jahre zu Gesicht bekommen und daher noch keine Rechnungen über dieselbe anstellen können.

\par\indent Zu der Anzeige meines Programms im Februar-Heft [worin von einem nicht genannten Verfasser ,,in dieser Zeitschrift die hauptsächlichsten Formeln, auf denen die Methode beruht kurz dargestellt“ und die Numerirungen der Formeln und Bezeichnungsweise der \textsc{Gauss}schen Schrift, Methodus peculiaris elevationem poli determinandi, Seite 39 u.\,f. dieses Bandes ungeändert aufgenommen worden sind] muss ich noch bemerken, dass dort irriger Weise (S. 136) von gleichen Höhen die Rede ist; eine solche Bedingung ist der Aufgabe ganz fremd. Ferner verdient noch folgender Zusatz angehängt zu werden: ,,Bei der Bestimmung von $\lambda$ durch die Tangente nach der Formel 9, bleibt die Zweideutigkeit zurück, ob man $\lambda$ im ersten oder zweiten Halbkreise nehmen müsse; diese wird durch Formel 4 gehoben, welche zeigt, dass $\lambda$ in demselben Halbkreise genommen werden muss, in welchem $u$ liegt, weil offenbar $\cos\varphi$ und $\cos h$ ihrer Natur nach positive Grössen sind.“

\par\indent Es versteht sich übrigens, dass dieser Zusatz blos der analytischen Vollständigkeit wegen noch nöthig war, denn man würde doch nicht zweifelhaft sein können, welcher Werth von $\lambda$ der rechte sei, da der zweite Werth den Stand der Uhr um 12 Stunden falsch und die Polhöhe mit dem unrechten Zeichen geben würde. Auch aus Ihrem Auszuge wird jeder sich die geometrischen Erörterungen, welche das Programm selbst enthält, leicht ergänzen und die Bedeutung der gebrauchten Zeichen auf der Sphäre nachweisen können. Übrigens hat Herr Prof. \textsc{Harding} für das astronomische Jahrbuch eine deutsche Übersetzung davon veranstaltet.
